\expandafter\def\csname HistoricalNoteText14\endcsname{Es ist leicht zu zeigen, dass auch für den Werth $r=R$ selbst eine gleiche Folgeordnung noch
gültig bleibt, nur mit der Einschränkung, dass dann in ganz speciellen Fällen ein Uebergangswerth von $\rho$,
(d.\,i. ein solcher, für welchen $T$ oder $U=0$ wird) mit einer der Grössen $-\omega,\ \omega,\ 3\omega,\ 5\omega$ u.\,s.\,w. zusammenfallen kann, während für alle grösseren Werthe von $r$ jeder Uebergangswerth von $\rho$ \emph{zwischen} zweien dieser Grössen liegen muss. Ich halte mich jedoch dabei nicht auf, da für unsern Zweck zureicht, das Bestehen
jener Folgeordnung, von irgend einem Werthe von $r$ an, nachgewiesen zu haben.}
